\subsection{Topologies compatible with a duality}

In this section, E and F denote two vector spaces put into duality by a bilinear form B (II, p.~40). We recall (II, p.~41) that we defined two linear mappings

\[
d _ {\mathbf {B}}: \mathrm{F} \to \mathrm{E} ^ {*}, s _ {\mathbf {B}}: \mathrm{E} \to \mathrm{F} ^ {*}
\]

characterized by the relation

\[
\mathbf {B} (x, y) = \langle x, d _ {\mathbf {B}} (y) \rangle = \langle y, s _ {\mathbf {B}} (x) \rangle\tag{1}
\]

for \(x\in \mathbf{E},y\in \mathbf{F}.\)

DEFINITION 1. --- A locally convex topology T on E is said to be compatible with the duality between E and F if \(d_{B}\) is a bijection from F onto the dual of the locally convex space obtained by assigning the topology T to E.

If there exists one such topology \(\mathcal{T}\) , the mapping \(d_{\mathrm{B}}\) is injective, that is to say, the duality between E and F is separating in F (II, p.~41).

PROPOSITION 1. --- (i) The closed convex subsets in E are the same for all the locally convex topologies on E which are compatible with the duality between E and F.

\begin{enumerate}
\def\labelenumi{(\roman{enumi})}
\setcounter{enumi}{1}
\tightlist
\item
  The bounded subsets of \(\mathbf{E}\) are the same for all the locally convex topologies on \(\mathbf{E}\) which are compatible with the duality between \(\mathbf{E}\) and \(\mathbf{F}\) .
\end{enumerate}

Let T be a topology on E compatible with the duality between E and F, hence finer than \(\sigma(\mathrm{E},\mathrm{F})\) . If a convex subset of E is closed for T, it is the intersection of closed, real half-spaces (II, p.~38, cor. 1), hence it is closed for \(\sigma(\mathrm{E},\mathrm{F})\) . This proves (i). Assertion (ii) was proved in cor. 3 of III, p.~27.

Let \(F_{\sigma}\) denote the vector space F endowed with the weak topology \(\sigma(F, E)\) . Then the linear mapping \(s_{B}\) maps E onto the dual \((F_{\sigma})'\) of \(F_{\sigma}\) (II, p.~43, prop. 3). Let \(S\) be a family of bounded subsets of \(F_{\sigma}\) . By abuse of language, the inverse image under \(s_{B}\) of the \(S\) -topology on \((F_{\sigma})'\) is called the \(S\) -topology on E. It is defined by the family of semi-norms

\[
p _ {\mathbf {A}} (x) = \sup _ {y \in \mathbf {A}} | \mathrm{B} (x, y) |,\tag{2}
\]

where A runs through S. In particular, when S is the family of finite subsets of F, the S-topology is precisely the weak topology \(\sigma(\mathrm{E}, \mathrm{F})\) .

DEFINITION 2. --- Let E and F be two spaces in duality. The Mackey topology on E, denoted by \(\tau(\mathrm{E},\mathrm{F})\) is defined as the \(\mathfrak{S}\) -topology on E, where \(\mathfrak{S}\) is the family of all subsets of F whose image in \(E^{*}\) (under \(d_{\mathrm{B}}\) ) is convex, balanced and compact for \(\sigma(\mathrm{E}^{*},\mathrm{E})\) .

When the duality between E and F is separating in F, \(d_{B}\) is injective and the topology \(\sigma(F, E)\) on F is the inverse image under \(d_{B}\) of the topology \(\sigma(E^{*}, E)\) on \(E^{*}\) . In this case, S consists of all those subsets of F which are convex, balanced and compact for \(\sigma(F, E)\) .

In general, if \(F_{1}=d_{\mathrm{B}}(\mathrm{F})\subset\mathrm{E}^{*}\) , and if we denote by \((x,y_{1})\mapsto\mathrm{B}_{1}(x,y_{1})\) the restriction of the canonical bilinear form \((x,x^{*})\mapsto\langle x,x^{*}\rangle\) to \(E\times F_{1}\) , then E and \(F_{1}\) are put in duality by \(B_{1}\) , and this duality is separating in \(F_{1}\) , since by definition we have \(\mathrm{B}(x,y)=B_{1}(x,d_{\mathrm{B}}(y))\) , def. 2 shows that \(\tau(\mathrm{E},\mathrm{F})=\tau(\mathrm{E},\mathrm{F}_{1})\) .

Remark 1. --- Let A be a compact convex subset of a Hausdorff locally convex space G, and let \(\tilde{A}\) be the closed convex balanced envelope of A. When the field K is R, the set \(\tilde{A}\) is the closed convex envelope of \(A \cup (-A)\) ; when K is C, the set \(\tilde{A}\) is contained in the closed convex envelope of \(2A \cup (-2A) \cup (2iA) \cup (-2iA)\) . Consequently (II, p.~14, prop. 15), \(\tilde{A}\) is compact.

We deduce, in particular, that when the duality between E and F is separating in F, the Mackey topology \(\tau(\mathrm{E},\mathrm{F})\) is also the \(S'\) -topology, where \(S'\) is the set of all convex subsets of F which are compact for \(\sigma(\mathrm{F},\mathrm{E})\) . In an analogous way we define the Mackey topology \(\tau(\mathrm{F},\mathrm{E})\) on F.

THEOREM 1 (Mackey). --- Let E and F be two spaces in duality; suppose that the duality is separating in F. In order that a locally convex topology T on E be compatible with the duality between E and F, it is necessary and sufficient that T be finer than the topology \(\sigma(\mathrm{E},\mathrm{F})\) and coarser than the Mackey topology \(\tau(\mathrm{E},\mathrm{F})\) .

Identify F with its image in \(E^{*}\) under \(d_{B}\) . Let \(S_{0}\) denote the set of all subsets of F which are convex, balanced and compact for \(\sigma(F, E)\) . By definition, \(\tau(E, F)\) is the \(S_{0}\) -topology on E, hence is finer than \(\sigma(E, F)\) .

Lemma 1. --- The subspace F of E* consists of all linear forms on E which are continuous for \(\tau(\mathrm{E},\mathrm{F})\) .

Every element of F is a continuous mapping for \(\sigma(\mathrm{E},\mathrm{F})\) , hence for \(\tau(\mathrm{E},\mathrm{F})\) .

Conversely, let \(f \in \mathrm{E}^*\) be continuous for \(\tau(\mathrm{E}, \mathrm{F})\) . There exists a neighbourhood \(\mathbf{U}\) of 0 in \(\mathrm{E}\) (for \(\tau(\mathrm{E}, \mathrm{F})\) ), such that \(|f| \leqslant 1\) on \(\mathbf{U}\) ; we can assume that there exists a set \(\mathbf{A} \in \mathfrak{S}_0\) such that \(\mathbf{U} = \mathbf{A}^\circ\) . In other words, \(f\) belongs to the bipolar \(\mathbf{A}^{\circ\circ}\) of \(\mathbf{A}\) for the duality between \(\mathbf{E}^*\) and \(\mathbf{E}\) . But the topology \(\sigma(\mathrm{F}, \mathrm{E})\) on \(\mathbf{F}\) is induced by \(\sigma(\mathrm{E}^*, \mathrm{E})\) ; consequently \(\mathbf{A}\) is convex, balanced and compact for \(\sigma(\mathrm{E}^*, \mathrm{E})\) , and the theorem of bipolars (II, p.~44, th. 1) implies the equality \(\mathbf{A} = \mathbf{A}^{\circ\circ}\) . Therefore we have that \(f \in \mathrm{F}\) , from which the lemma follows.

Lemma 2. --- Let T be a locally convex topology on E such that every linear form on E which is continuous for T belongs to F. Then T is coarser than \(\tau(\mathrm{E}, \mathrm{F})\) .

Let \(\mathfrak{U}\) be the set of convex, balanced neighbourhoods of 0 for \(\mathcal{T}\) . Let \(\mathfrak{S}\) be the set of polars in F of elements of \(\mathfrak{U}\) . By cor. 2 of III, p.~17, we have \(\mathfrak{S} \subset \mathfrak{S}_0\) , and by cor. 1 of prop. 7 of III, p.~19, \(\mathcal{T}\) is identical with the \(\mathfrak{S}'\) -topology, where \(\mathfrak{S}'\) is the set of polars of sets of \(\mathfrak{U}\) in the dual E' of E. But E' \(\subset\) F, by hypothesis, hence every set of \(\mathfrak{S}'\) is contained in a set of \(\mathfrak{S}\) ; and the lemma follows.

Let T be a topology on E compatible with the duality between E and F. Then T is coarser than \(\tau(\mathrm{E},\mathrm{F})\) by lemma 2, and evidently T is finer than \(\sigma(\mathrm{E},\mathrm{F})\) . Conversely, F is the dual of E for the topology \(\tau(\mathrm{E},\mathrm{F})\) (lemma 1) and for the topology \(\sigma(\mathrm{E},\mathrm{F})\) (II, p.~43, prop. 3), hence also for every topology intermediate between \(\tau(\mathrm{E},\mathrm{F})\) and \(\sigma(\mathrm{E},\mathrm{F})\) .

COROLLARY. --- Let p be a semi-norm on E. The following conditions are equivalent :

\begin{enumerate}
\def\labelenumi{(\roman{enumi})}
\item
  \(p\) is continuous for the topology \(\tau(\mathrm{E},\mathrm{F})\) ;
\item
  every linear form \(f\) on \(\mathbf{E}\) , such that \(|f| \leqslant p\) , comes from an element of \(\mathbf{F}\) .
\item[]
  \((i) \Rightarrow (ii):\) if \(p\) is continuous for \(\tau(\mathrm{E},\mathrm{F})\) , every linear form \(f\) on \(\mathrm{E}\) such that \(|f|\leqslant p\) is continuous for \(\tau(\mathrm{E},\mathrm{F})\) , hence comes from an element of \(\mathrm{F}\) by lemma 1.
\item[]
  \((ii) \Rightarrow (i):\) let T be the topology on E defined by the semi-norm p.~If condition (ii) is satisfied, the linear forms on E which are continuous for T belong to F. By lemma 2 T is coarser than \(\tau(\mathrm{E}, \mathrm{F})\) , hence p is continuous for \(\tau(\mathrm{E}, \mathrm{F})\) .
\end{enumerate}

Remark 2. --- * Let K be a convex subset of F which is compact for the weak topology \(\sigma(F, E)\) and \(\mu\) a positive measure on K. Put

\[
p (x) = \int_ {\mathbf {K}} | \mathrm{B} (x, y) | d \mu (y)
\]

for all \(x \in E\) . It is immediate that p is a semi-norm. Moreover, for every \(x \in E\) , the relation \(\langle |B(x, y)| \leqslant 1 \text{ for all } y \in K \rangle\) implies that \(p(x) \leqslant \mu(K)\) . This proves that the semi-norm p on E is continuous for the Mackey topology \(\tau(E, F)\) . *

Example. --- Let G be a locally convex space and \(G'\) its dual. On \(G'\) , the weak topology \(\sigma(G', G)\) and the topology of convex compact convergence (III, p.~14) are compatible with the duality between \(G'\) and G. In general, the strong topology and the topology of compact convergence on \(G'\) are not compatible with the duality between \(G'\) and G. Recall however that when G is Hausdorff and quasi-complete, the topology of compact convergence on \(G'\) coincides with that of convex compact convergence (III, p.~8), hence is compatible with the duality between \(G'\) and G.

DEFINITION 3. --- Let E and F be two vector spaces in duality, and T the family of subsets of F which are bounded for \(\sigma(\mathrm{F}, \mathrm{E})\) . Then the T-topology on F is denoted by \(\beta(\mathrm{E}, \mathrm{F})\) .

Similarly, we define the topology \(\beta(\mathrm{F},\mathrm{E})\) on F. It can be seen easily that the topology \(\beta(\mathrm{E},\mathrm{F})\) is identical with \(\beta(\mathrm{E},\mathrm{F}/\mathrm{E}^{\circ})\) , and we can reduce to the case when the duality between E and F is separating in F.

Remarks. --- 3) Let \(E_{\sigma}\) denote the space E endowed with the topology \(\sigma(\mathrm{E}, \mathrm{F})\) . The barrels (III, p.~24) in \(E_{\sigma}\) are the subsets of E which are convex, balanced closed and absorbent for \(\sigma(\mathrm{E}, \mathrm{F})\) . These are none other than the polars of the subsets of F which are convex, balanced and bounded for \(\sigma(\mathrm{F}, \mathrm{E})\) . Consequently, the family of all barrels in \(E_{\sigma}\) is a fundamental system of neighbourhoods of 0 for the topology \(\beta(\mathrm{E}, \mathrm{F})\) in E. In other words, a semi-norm on E is continuous for \(\beta(\mathrm{E}, \mathrm{F})\) if and only if it is lower semi-continuous for \(\sigma(\mathrm{E}, \mathrm{F})\) (cf.~III, p.~24, prop. 1).

\begin{enumerate}
\def\labelenumi{\arabic{enumi})}
\setcounter{enumi}{3}
\item
  Let \(\mathcal{T}\) be a topology on E compatible with the duality between E and F. By prop. 1, (ii) of IV, p.~1, the topology \(\beta(\mathrm{F},\mathrm{E})\) on F is none other than the strong topology on F, when F is identified with the dual of E (with the topology \(\mathcal{T}\) ).
\item
  The topology \(\beta(\mathrm{E},\mathrm{F})\) on \(\mathrm{E}\) is finer than \(\tau(\mathrm{E},\mathrm{F})\) . It is not, in general compatible with the duality between \(\mathrm{E}\) and \(\mathrm{F}\) (cf.~however § 2). In particular, a subset of \(\mathrm{E}\) which is bounded for \(\sigma(\mathrm{E},\mathrm{F})\) is not necessarily bounded for \(\beta(\mathrm{E},\mathrm{F})\) .
\end{enumerate}
% semantic-fragments: legacy-input-seam
\relax{}
